% Generated by scripts/check_manual_examples.py --write; edit that source.
\section{Verified runnable examples}
\label{chap:current-examples}
These examples are checked against the current checkout by \texttt{scripts/check\_manual\_examples.py}. Run it with Python and the core package dependencies installed; no network, credentials, or optional scientific extras are needed. Each listing is a standalone Python program. The output listings are exact stdout checked by the script, including the explicit formatting of numerical results. Exception examples deliberately display only the exception class, not unstable full error messages.

\begin{lstlisting}[style=mikoshi]
python scripts/check_manual_examples.py
python scripts/check_manual_examples.py --write
\end{lstlisting}
Without \texttt{--write}, the script also checks that this generated chapter has not drifted. With \texttt{--write}, it regenerates the chapter only after all examples pass. This is a bounded documentation check, not a claim to verify the whole library.

\subsection{Session: Python data and exact arithmetic}
Use \texttt{Session.eval(source)} for Python data. Exact rationals become \texttt{fractions.Fraction}; an inexact zero remains a Python float. \texttt{Session.evaluate(source)} instead returns the raw language result.

\begin{lstlisting}[style=mikoshi]
from fractions import Fraction
from mikoshilang import Session

s = Session()
assert s.eval("1 + 2") == 3
q = s.eval("1/3")
assert type(q) is Fraction and q == Fraction(1, 3)
z = s.eval("0.5 - 0.5")
assert type(z) is float and z == 0.0
print(s.eval("1 + 2"))
print(repr(q))
print(repr(z))
\end{lstlisting}
\noindent Expected standard output:
\begin{lstlisting}[style=mikoshi]
3
Fraction(1, 3)
0.0
\end{lstlisting}

\subsection{Session: definitions, patterns, and isolation}
Definitions and pattern dispatch belong to one session. The integer pattern is more specific than the general pattern, and the literal definition is more specific still.

\begin{lstlisting}[style=mikoshi]
from mikoshilang import Session

s = Session()
s.eval('f[x_] := "general"; f[n_Integer] := n^2; f[0] = 7')
values = s.eval(' {f[0], f[3], f["x"]} ')
assert values == [7, 9, "general"]
assert str(Session().eval("f[3]")) == "f[3]"
print(values)
print(Session().eval("f[3]"))
\end{lstlisting}
\noindent Expected standard output:
\begin{lstlisting}[style=mikoshi]
[7, 9, 'general']
f[3]
\end{lstlisting}

\subsection{Session: solution rules, ODEs, and series}
The language heads \texttt{Solve} and \texttt{DSolve} return nested rule lists, not bare roots or Python dictionaries. Here \texttt{evaluate} preserves the raw language braces; \texttt{eval} converts the list containers to Python lists while leaving rules as expressions. Other symbolic results remain expressions by default. Series includes the requested order; \texttt{Normal} removes the remainder term.

\begin{lstlisting}[style=mikoshi]
from mikoshilang import Session

s = Session()
print(s.evaluate("Solve[{x + y == 3, x - y == 1}, {x, y}]"))
print(s.evaluate("DSolve[{y'[x] == y[x], y[0] == 1}, y[x], x]"))
print(s.eval("Series[Exp[x], {x, 0, 3}]"))
print(s.eval("Normal[Series[Exp[x], {x, 0, 3}]]"))
assert s.eval("{x, y} /. First[Solve[{x+y==3, x-y==1}, {x,y}]]") == [2, 1]
assert s.eval("Normal[Series[Exp[x], {x,0,3}]] /. x -> 0") == 1
\end{lstlisting}
\noindent Expected standard output:
\begin{lstlisting}[style=mikoshi]
{{x -> 2, y -> 1}}
{{y[x] -> E^x}}
1 + x + x^2/2 + x^3/6 + O[x]^4
1 + x + x^2/2 + x^3/6
\end{lstlisting}

\subsection{Session: verification status and strict type predicates}
Only the outer status is displayed here, not the full evidence. \texttt{Certified} means an exact certificate was checked; \texttt{Checked} means symbolic or numerical evidence, not a proof; \texttt{Unchecked} means no applicable or conclusive checker. A failed check returns \texttt{Failed}; computation errors still propagate. \texttt{TypeMatchQ} is a strict predicate: a symbolic value is not known to be an integer. This differs from opt-in call-boundary type checking, which permits symbolic arguments.

\begin{lstlisting}[style=mikoshi]
from mikoshilang import Session

s = Session()
print(s.evaluate("Head[Verify[GCD[12, 18]]]"))
print(s.evaluate("Head[Verify[Integrate[x^2, x]]]"))
print(s.evaluate("Head[Verify[Sort[{3, 1, 2}]]]"))
matches = s.eval("{TypeMatchQ[{1,2}, ListOf[Integer]], TypeMatchQ[x, Integer]}")
assert matches == [True, False]
print(matches)
\end{lstlisting}
\noindent Expected standard output:
\begin{lstlisting}[style=mikoshi]
Certified
Checked
Unchecked
[True, False]
\end{lstlisting}

\subsection{Session: expose a typed Python function locally}
The host deliberately exposes a function with \texttt{@s.function}; this needs no Python-import permission and does not register a global library head. Type annotations constrain concrete inputs. The error example prints only the checked exception class.

\begin{lstlisting}[style=mikoshi]
from mikoshilang import Session

s = Session()

@s.function(listable=True, check="exact")
def square_plus_one(n: int) -> int:
    return n * n + 1

values = s.eval("SquarePlusOne[{1, 2, 3}]")
assert values == [2, 5, 10]
assert str(Session().eval("SquarePlusOne[2]")) == "SquarePlusOne[2]"
print(values)
try:
    s.eval("SquarePlusOne[2.5]")
except TypeError as exc:
    print(type(exc).__name__)
else:
    raise AssertionError("expected a TypeError")
\end{lstlisting}
\noindent Expected standard output:
\begin{lstlisting}[style=mikoshi]
[2, 5, 10]
TypeError
\end{lstlisting}

\subsection{Session: explicit permission for Python calls}
Python imports/calls are disabled by default. An allow-list enables the named modules and their submodules; it is not a security sandbox. Do not execute untrusted input without an external process sandbox. This offline example uses only the standard library.

\begin{lstlisting}[style=mikoshi]
from mikoshilang import Session

call = 'PythonFunction["math.gcd"][12, 18]'
try:
    Session().eval(call)
except PermissionError as exc:
    print(type(exc).__name__)
else:
    raise AssertionError("Python calls must be gated")
s = Session(allow_python={"math"})
assert s.eval(call) == 6
print(s.eval(call))
\end{lstlisting}
\noindent Expected standard output:
\begin{lstlisting}[style=mikoshi]
PermissionError
6
\end{lstlisting}

\subsection{Direct Python: List, Take, Drop, and Part}
These imported helpers are not \texttt{Session} language calls. \texttt{Take(lst, n)}, \texttt{Drop(lst, n)}, and \texttt{Part(lst, *indices)} accept a \texttt{List} expression, not an ordinary Python list. Negative counts select/drop from the end; zero counts select nothing/drop nothing. Part indices are one-based, negative indices count from the end, and index zero returns the head.

\begin{lstlisting}[style=mikoshi]
from mikoshilang import List, Take, Drop, Part

xs = List(10, 20, 30, 40)
assert Take(xs, -2).args == (30, 40)
assert Drop(xs, -2).args == (10, 20)
assert Take(xs, 0).args == ()
assert Drop(xs, 0) == xs
assert Part(xs, -1) == 40
print(Take(xs, -2))
print(Drop(xs, -2))
print(Take(xs, 0))
print(Drop(xs, 0))
print(Part(xs, 0))
print(Part(xs, -1))
try:
    Take([10, 20], 1)
except TypeError as exc:
    print(type(exc).__name__)
else:
    raise AssertionError("expected a List expression requirement")
\end{lstlisting}
\noindent Expected standard output:
\begin{lstlisting}[style=mikoshi]
{30, 40}
{10, 20}
{}
{10, 20, 30, 40}
List
40
TypeError
\end{lstlisting}

\subsection{Direct Python: finite Range and callable Table}
\texttt{Range(n)}, \texttt{Range(a, b)}, and \texttt{Range(a, b, step)} use integer bounds with an inclusive end. Direct \texttt{Table(fn, spec)} takes a Python callable and either a count, \texttt{(min, max)}, or \texttt{(min, max, step)}; it does not take a language expression and iterator symbol. Wrong-direction steps return an empty list rather than looping. Table rejects zero, nonfinite, and nonadvancing steps. The verifier runs this entire example in a timeout-limited child process; the checks below include wrong-direction and nonadvancing cases.

\begin{lstlisting}[style=mikoshi]
from mikoshilang import Range, Table

assert Range(3, 1, -1).args == (3, 2, 1)
assert Range(1, 3, -1).args == ()
assert Table(lambda i: i*i, (3, 1, -1)).args == (9, 4, 1)
assert Table(lambda i: i, (1, 3, -1)).args == ()
print(Range(3, 1, -1))
print(Range(1, 3, -1))
print(Table(lambda i: i*i, (3, 1, -1)))
print(Table(lambda i: i, (1, 3, -1)))
for spec in [(1, 3, 0), (1, 3, float("inf")),
             (1e20, 1e20 + 1e6, 1.0)]:
    try:
        Table(lambda i: i, spec)
    except ValueError as exc:
        print(type(exc).__name__)
    else:
        raise AssertionError("expected invalid-step rejection")
\end{lstlisting}
\noindent Expected standard output:
\begin{lstlisting}[style=mikoshi]
{3, 2, 1}
{}
{9, 4, 1}
{}
ValueError
ValueError
ValueError
\end{lstlisting}

\subsection{Direct Python: tiny determinants and complex eigenvalues}
Direct linear-algebra helpers use \texttt{Matrix(*rows)}, whereas Session language calls use nested lists. Integer determinant input is exact; floating input is numerical and a tiny nonzero determinant must not be rounded to zero. Complex eigenvalues keep their imaginary parts. Formatting below normalizes insignificant numerical variation; the assertions check values, not just display or shape.

\begin{lstlisting}[style=mikoshi]
from math import isclose
from mikoshilang import Matrix, Det, Eigenvalues

assert Det(Matrix([2, 0], [0, 3])) == 6
tiny = Det(Matrix([1e-12, 0.0], [0.0, 1.0]))
assert tiny != 0 and isclose(tiny, 1e-12, rel_tol=1e-12, abs_tol=0)
values = Eigenvalues(Matrix([1j, 0], [0, -1j])).args
assert len(values) == 2 and set(values) == {-1j, 1j}
print(f"{tiny:.0e}")
print([complex(v) == -1j or complex(v) == 1j for v in values])
print([v.imag for v in values])
\end{lstlisting}
\noindent Expected standard output:
\begin{lstlisting}[style=mikoshi]
1e-12
[True, True]
[-1.0, 1.0]
\end{lstlisting}

\subsection{Direct Python: units, balancing, and complex convolution}
\texttt{UnitConvert(Quantity(value, unit), target)} returns a Python Quantity object. Offset-temperature conversion is supported, but offset-temperature arithmetic is rejected: convert to kelvin first. \texttt{BalanceEquation(equation)} accepts an equation string; \texttt{Convolve(a, b)} accepts List expressions and preserves complex components. These examples require no SciPy extra.

\begin{lstlisting}[style=mikoshi]
from mikoshilang import (
    Quantity, UnitConvert, BalanceEquation, Convolve, List,
)

distance = UnitConvert(Quantity(1, "km"), "m")
assert distance.value == 1000.0 and distance.unit_name == "m"
temperature = UnitConvert(Quantity(0, "C"), "K")
assert temperature.value == 273.15 and temperature.unit_name == "K"
print(distance)
print(temperature)
try:
    Quantity(0, "C") + Quantity(1, "C")
except ValueError as exc:
    print(type(exc).__name__)
else:
    raise AssertionError("offset-temperature arithmetic must be rejected")
balanced = BalanceEquation("H2 + O2 -> H2O")
assert balanced == "2H2 + O2 -> 2H2O"
print(balanced)
result = Convolve(List(1j, 1), List(1, 1j))
assert result.args == (1j, 0, 1j)
print(result.args)
\end{lstlisting}
\noindent Expected standard output:
\begin{lstlisting}[style=mikoshi]
Quantity[1000, "m"]
Quantity[273.15, "K"]
ValueError
2H2 + O2 -> 2H2O
(1j, 0, 1j)
\end{lstlisting}
